\documentclass[legal]{polinetwork}

\pnTitle{Informativa Privacy}
\pnSubtitle{Admin e Volontari}

\begin{document}

\pnMakeHeader

\section{Titolare del trattamento}

Ai sensi degli artt. 4 e 24 del Reg. UE 2016/679 il titolare del trattamento è individuabile nella figura del legale rappresentante dell’associazione PoliNetwork APS, il suo Presidente.

L’associazione non ha nominato un Responsabile della Protezione dei Dati (DPO) in quanto non ricorrono le condizioni di cui all’art. 37 del GDPR.

\textbf{Contatti per esercitare i diritti privacy:} \\

\textbf{Indirizzo email:} \href{mailto:privacy@polinetwork.org}{privacy@polinetwork.org} (in alternativa \href{mailto:direttivo@polinetwork.org}{direttivo@polinetwork.org}) \\
\textbf{PEC:} \href{mailto:polinetwork@pec.it}{polinetwork@pec.it}

\section{Tipologia di dati raccolti}

Il Titolare del Trattamento tratta i dati personali identificativi comunicati dall’interessato in occasione della richiesta di volontariato presso PoliNetwork APS:

\begin{itemize}
	\item Nome; Cognome; Eventuale Alias;
	\item Email; Numero di telefono;
	\item Username impostato sulla piattaforma Telegram;
	\item Data di inizio attività di volontariato;
	\item Informazioni sugli studi, sull’eventuale attività di rappresentanza presso il Politecnico di Milano e ruoli all’interno dell’associazione e/o in altre associazioni;
	\item Eventuali informazioni aggiuntive fornite dal volontario;
	\item Codice Fiscale (richiesto per i membri del team HR).
\end{itemize}

\section{Finalità e base giuridica del trattamento}

I dati di natura personale forniti saranno oggetto di trattamento nel rispetto delle condizioni di liceità previste dal Reg. UE 2016/679 (in particolare ex art. 6 lett. b, lett. c e lett. f), per l’instaurazione, la gestione organizzativa e la pianificazione del rapporto di collaborazione come volontario per l’associazione, ed in particolare per:

\begin{itemize}
	\item inserimento nelle anagrafiche interne, assegnazione dei ruoli operativi e gestione dei relativi permessi o accessi;
	\item coordinamento delle attività, gestione delle mansioni e comunicazione interna necessaria allo svolgimento del ruolo;
	\item archiviazione e analisi organizzativa interna, tra cui il monitoraggio della composizione dei gruppi di lavoro e la valutazione della disponibilità dei collaboratori (ad esempio in vista della conclusione del percorso di studi), al fine di pianificare e avviare tempestivamente nuove campagne di recruitment;
	\item adempiere agli obblighi contrattuali, di legge e finalità amministrativo-contabile. Ai fini dell’applicazione delle disposizioni in materia di protezione dei dati personali, i trattamenti effettuati per finalità amministrativo-contabili sono quelli connessi allo svolgimento delle attività di natura organizzativa, amministrativa e gestionale dei collaboratori;
	\item adempiere agli obblighi previsti dalla legge (come ad esempio la stipula delle coperture assicurative obbligatorie per i volontari e la tenuta dei registri del Terzo Settore), da un regolamento, dalla normativa comunitaria o da un ordine dell’Autorità;
	\item esercitare i diritti del Titolare, ad esempio il diritto di difesa in giudizio.
\end{itemize}

Il conferimento dei dati personali per le finalità sopra citate è necessario per dare seguito all’instaurazione del rapporto di collaborazione. Il mancato conferimento dei dati personali può comportare l’impossibilità di assumere il ruolo di amministratore o volontario all'interno dell'associazione.

\section{Modalità del trattamento}

Il trattamento dei dati personali avviene mediante strumenti manuali, informatici e telematici. Le logiche applicate sono strettamente correlate alle finalità stesse e implementate in modo da garantire la massima sicurezza e riservatezza dei dati, prevenendone la perdita, gli usi illeciti e gli accessi non autorizzati.

\section{Conservazione dei dati}

Il trattamento sarà svolto in forma automatizzata e/o manuale, con modalità e strumenti volti a garantire la massima sicurezza e riservatezza, ad opera di soggetti a ciò appositamente incaricati. Nel rispetto di quanto previsto dall’art. 5 comma 1 lett. e) del Reg. UE 2016/679 i dati personali raccolti verranno conservati in una forma che consenta l’identificazione degli interessati per un arco di tempo non superiore al conseguimento delle finalità per le quali i dati personali sono trattati.

\section{Destinatari e accesso ai dati}

I dati personali non saranno comunicati a soggetti terzi esterni all’associazione né diffusi pubblicamente. I dati potranno essere accessibili esclusivamente ai membri dell’associazione PoliNetwork APS direttamente coinvolti nella gestione amministrativa, organizzativa e operativa dei volontari (ad esempio, membri del Consiglio Direttivo, referenti dell'area HR o coordinatori dei progetti in cui il volontario è coinvolto). Tali soggetti sono stati appositamente istruiti e tratteranno i dati nel pieno rispetto delle direttive interne di riservatezza.

Potranno inoltre avere accesso ai dati i soggetti che forniscono servizi per la gestione e la manutenzione dei sistemi informatici del Titolare, i quali operano in qualità di Responsabili del trattamento ai sensi dell’art. 28 del GDPR. In particolare, per la raccolta dei dati e delle adesioni dei volontari viene utilizzata la piattaforma Microsoft Forms, fornita da Microsoft Corporation.

\textbf{Trasferimento dei dati verso Paesi terzi:} Microsoft Corporation può trattare dati personali negli Stati Uniti d’America. Tale trasferimento avviene nel rispetto del GDPR sulla base delle seguenti garanzie:

\begin{itemize}
	\item l’adesione di Microsoft al \textbf{Data Privacy Framework (DPF) UE-USA}, adottato dalla Commissione Europea con decisione di adeguatezza del 10 luglio 2023 (ai sensi dell’art. 45 del GDPR). L’elenco delle organizzazioni aderenti al DPF è consultabile al sito;
	\item la stipula di \textbf{Clausole Contrattuali Standard (SCC)} approvate dalla Commissione Europea (ai sensi dell’art. 46 del GDPR), che Microsoft incorpora nelle proprie condizioni contrattuali come misura di garanzia complementare.
\end{itemize}

Per ulteriori informazioni sulle modalità di trattamento dei dati da parte di Microsoft si rimanda alla relativa informativa privacy, consultabile all’indirizzo \href{https://privacy.microsoft.com}{privacy.microsoft.com}.

\section{Diritti degli interessati}

L’interessato potrà far valere i propri diritti come espressi dagli artt. 15-22 del Regolamento UE 2016/679, rivolgendosi al Titolare del Trattamento tramite i recapiti indicati nella Sezione 1 (”Titolare del trattamento”) del presente documento.

Nello specifico, l’interessato ha il diritto, in qualunque momento, di:
\begin{itemize}
	\item \textbf{Diritto di accesso (art. 15):} ottenere la conferma dell’esistenza o meno di dati personali che lo riguardano, accedervi e riceverne copia;
	\item \textbf{Diritto di rettifica (art. 16):} richiedere l’aggiornamento, la rettifica di dati inesatti o l’integrazione di quelli incompleti;
	\item \textbf{Diritto alla cancellazione o ”all’oblio” (art. 17):} ottenere la cancellazione dei propri dati qualora non siano più necessari per le finalità per le quali sono stati raccolti, o in caso di revoca del consenso;
	\item \textbf{Diritto di limitazione del trattamento (art. 18):} richiedere che i dati siano unicamente conservati (e non altrimenti trattati) in determinate ipotesi previste dalla legge;
	\item \textbf{Diritto alla portabilità dei dati (art. 20):} ricevere in un formato strutturato, di uso comune e leggibile da dispositivo automatico i dati personali forniti e, se tecnicamente fattibile, trasmetterli a un altro titolare senza impedimenti;
	\item \textbf{Diritto di opposizione (art. 21):} opporsi in tutto o in parte, per motivi connessi alla propria situazione particolare, al trattamento dei dati personali che lo riguardano.
\end{itemize}

Si informa inoltre che il processo di selezione non è soggetto a decisioni basate unicamente sul trattamento automatizzato, ivi compresa la profilazione, che producano effetti giuridici o che incidano significativamente sulla persona (art. 22 del GDPR). Le decisioni relative alle candidature sono sempre adottate da persone fisiche.

Fatto salvo ogni altro ricorso amministrativo e giurisdizionale, se il candidato ritiene che il trattamento dei dati violi quanto previsto dal GDPR, ha il diritto di proporre reclamo al Garante per la protezione dei dati personali ai sensi dell’art. 77 del GDPR (\href{https://www.garanteprivacy.it}{www.garanteprivacy.it}).

Infine, conformemente all’art. 7, paragrafo 3, e all’art. 6, paragrafo 1, lettera a) del GDPR, l’interessato ha il diritto di revocare il proprio consenso in qualsiasi momento, senza che ciò pregiudichi la liceità del trattamento basata sul consenso prestato prima della revoca.

\end{document}
